\section{Introduction}

\label{sec:ray-introduction}

Existing metastable consensus protocols use fixed parameters $(k, \alpha, \beta)$ tuned
for worst-case adversarial conditions. This means that even in benign network conditions
(no Byzantine validators, low latency), transactions must wait for the same number of
confirmation rounds as under attack. This paper introduces the Ray protocol, which
adapts its finality threshold in real-time based on observed network behavior.

The key insight is that the \emph{rate} at which confidence accumulates carries
information about the adversarial environment. Rapid, unanimous convergence suggests
benign conditions where fewer confirmation rounds suffice. Slow, contested convergence
suggests adversarial interference requiring more rounds.

\paragraph{Contributions.}
\begin{enumerate}[nosep]
\item The Ray adaptive finality protocol with confidence velocity metric (\S\ref{sec:ray-protocol}).
\item An online Bayesian adversary estimator derived from query response distributions (\S\ref{sec:ray-estimator}).
\item A proof that adaptive thresholding preserves $2^{-128}$ safety (\S\ref{sec:ray-safety}).
\item A 40\% latency reduction compared to static protocols in benign conditions (\S\ref{sec:ray-evaluation}).
\end{enumerate}

%% -----------------------------------------------------------------------
%%  2. SYSTEM MODEL
%% -----------------------------------------------------------------------
